\begin{abstract}
When shared forcing synchronizes coupled dynamical systems, the
coupling itself becomes difficult to infer---an optimizer may converge
to a plausible-looking estimate that carries no statistical support.
We quantify this risk in the context of spatial epidemiology using the
Cram\'er-Rao Lower Bound (CRLB). For a discrete-time two-region SIR
model with seasonal forcing, we derive exact sensitivity equations and
compute the Fisher Information for the connectivity parameter
$\theta$. Identifiability depends on the degree of symmetry between
regions: when both seasonal phases and initial epidemic states are
similar, the CRLB diverges and $\theta$ cannot be reliably estimated.
Surprisingly, any source of asymmetry---different seasonal timing,
different initial conditions, or even different population
sizes---restores identifiability. We validate these predictions on
synthetic data and apply the framework to pneumonia and influenza
excess mortality across US states. Despite 12 flu seasons of data,
connectivity estimates for nearly all state pairs cannot be
distinguished from zero, consistent with the strong synchronization
of temperate influenza epidemics. For diseases where regional
asymmetries are naturally present, such as measles with heterogeneous
vaccination coverage, our analysis predicts that connectivity
inference from case data may be feasible.
\end{abstract}
